\chapter{Introduction}
\section{Motivation}

Energy-efficient and low-latency classification of neuromuscular activity is crucial in many medical applications. For example, assistive devices for people with impaired lower-limb mobility aim to support movements that can no longer be performed independently due to nerve injury. For such a device to be useful in everyday life, the system has to determine which movement the user intends to perform and react accordingly. Surface \ac{EMG} provides a way to obtain this information from the remaining voluntary muscle activity and can therefore serve as an input for intent recognition and the control of assistive devices \cite{cimolato_emg-driven_2022}.

The classification of \ac{EMG} signals, however, represents a challenging embedded computing problem. A wearable controller must process a continuous stream of sensor data and provide classifications with sufficiently low latency while operating with limited computational resources, memory, and energy. At the same time, classification accuracy and robustness cannot simply be sacrificed for efficiency, since incorrect or delayed predictions directly affect the behavior of the controlled device. Consequently, the classifier has to provide a suitable trade-off between classification performance and computational cost \cite{castruita-lopez_electromyography_2025}.\\

\ac{HDC} is a promising approach for this application. \ac{HDC} represents information using high-dimensional vectors and processes these representations using comparatively simple and highly parallel vector operations. Training and inference can therefore be implemented without the complex computational structures required by many conventional learning approaches, making HDC particularly interesting for real-time processing on resource-constrained embedded devices \cite{thomas_theoretical_2021}.

At the same time, the properties that make \ac{HDC} attractive also introduce considerable implementation cost. Hypervectors typically contain thousands of dimensions, and although the individual vector elements and operations are simple, every additional dimension increases the amount of data that has to be stored and processed. A similar trade-off exists also in the quantization and representation of continuous input values. Using a larger number of quantization levels can preserve more information about the input, but also increases the memory required to store these representations. The resulting implementation cost therefore depends strongly on parameters such as the vector dimension and the number of quantization levels. Reducing these parameters can therefore directly decrease memory consumption, the number of operations, and ultimately the hardware resources required by an implementation. However, reducing them too far can decrease the representational capacity of the model and lower its classification accuracy and robustness.

This problem motivates the central approach of this thesis. To reduce the computational and memory requirements without losing accuracy, the HDC model itself is first systematically optimized and adapted to the underlying EMG classification problem. The goal is to improve how information is represented and mapped to the hyperdimensional space such that relevant characteristics of the input data are preserved more effectively. By adapting the representation to the characteristics of the EMG signals, comparable or even improved classification performance may be achieved using smaller hypervectors. This creates the possibility of reducing the computational and memory requirements of HDC without necessarily sacrificing classification performance, making it more suitable for efficient implementation in resource-constrained embedded hardware.

\section{Goal}

The goal of this thesis is therefore to optimize an existing \ac{HDC} classifier for EMG-based foot movement recognition with respect to both classification performance and hardware efficiency. The objective is not only to maximize classification accuracy, but to determine how far the computational and memory requirements of the model can be reduced while maintaining or improving the accuracy of the original configuration.\\

A first focus is the vector dimension. The baseline \ac{HDC} model uses high-dimensional binary vectors, for which storage requirements and the computational cost of most vector operations scale directly with the number of dimensions. The influence of the vector dimension on classification accuracy is therefore investigated to determine how far it can be reduced without a relevant loss of classification performance.

Beyond reducing the vector dimension, this work investigates whether the available dimensions can be used more effectively. In the baseline model, continuous input values are mapped to a sequence of quantization levels whose corresponding hypervectors are separated by equal distances. This implies that equal differences in the input values are encoded into equal distances in the hyperdimensional space. For classification, however, some value ranges may contain more relevant information for distinguishing movement classes than others. A genetic algorithm is therefore used to optimize the distribution of distances between hypervectors representing consecutive quantization levels individually for each EMG feature. The objective is to improve class separability and classification accuracy and, in particular, to compensate for the loss of representational capacity that may result from reducing the vector dimension.\\

Besides the vector representation of quantization levels, a second optimization target affecting the size of the HDC model is the number of quantization levels. Reducing this number directly decreases the number of hypervectors that need to be stored, independently of the vector dimension. Different quantization strategies are therefore evaluated to determine whether the continuous EMG feature space can be represented using fewer, more informative levels than with conventional uniform quantization. In this way, both the placement of the quantization boundaries in the input domain and the arrangement of the corresponding hypervectors in hyperspace are optimized with the aim of preserving classification-relevant information more efficiently.\\

These algorithmic optimizations are complemented by implementation-level improvements that provide a hardware-oriented processing structure. Based on the optimized HDC model, a pipelined accelerator is implemented in SystemC and synthesized to RTL using High-Level Synthesis for an FPGA that is suitable for embedded applications. This step connects the algorithmic optimization to a concrete hardware implementation and enables the resulting resource requirements and processing performance to be evaluated.\\

Consequently, the work addresses four main questions:

\begin{enumerate}
\item Can the data representation in hyperdimensional space be optimized using a genetic algorithm to improve the classification accuracy or robustness of HDC? 

\item Can data-adaptive quantization methods further improve classification performance compared with conventional uniform quantization?

\item Based on these optimizations, how far can the vector dimension and the number of quantization levels be reduced while preserving classification accuracy?

\item Can the resulting HDC model be mapped efficiently to a parallel and pipelined FPGA accelerator, and what hardware cost and processing performance result from this implementation?

\end{enumerate}

Together, these investigations aim to derive an HDC configuration that does not use high dimensionality and large memory capacity simply by convention, but instead allocates them according to the requirements of the classification task. The resulting model and hardware implementation provide a basis for evaluating \ac{HDC} as an efficient classifier for embedded \ac{EMG} processing.